% Example CV

% Usar paquetes compartidos
\makeatletter
\edef\input@path{%
  {\@currdir../shared/moderncv/}%
  {\@currdir../shared/moderntimeline/}%
  {\@currdir../shared/pdfpages/}%
  {\@currdir../shared/xpatch/}%
  {\@currdir../shared/biblatex_modifications/}%
  {\@currdir../shared/cvitem_modifications/}%
  \input@path%
}
\makeatother

\documentclass[11pt, a4paper]{moderncv}
\moderncvcolor{blue}
\moderncvstyle{banking}
\usepackage[utf8]{inputenc}
\usepackage[T1]{fontenc}
\usepackage{lmodern}
\usepackage{graphicx}
\usepackage{tikz}
\usepackage[scale=0.8]{geometry}
\AtBeginDocument{\recomputelengths}

\firstname{Arturo}
\familyname{Veras Olivos}
\title{Curriculum Vitae}
\address{}{}
\mobile{+56 9 82413883}
\email{a.veras@gmail.com}

\begin{document}
\maketitle

\section{Profile}
Deep Tech Systems Architect and C++ Senior Engineer with nearly two decades of experience in hardware-resilient systems and algorithmic development. For the past 4 years, I have pioneered AI-Augmented Engineering workflows, architecting custom "Constitutions" (via dev-agents) that orchestrate AI swarms (Copilot, Cursor, LLMs) to generate MISRA-C compliant, ASan-verified, and fault-tolerant code. I do not just write code; I design robust architectures, enforce strict enterprise engineering standards, and multiply delivery velocity by 10x through advanced Prompt Engineering and automated QA pipelines.

\section{Experience}
\subsection{Professional Position}
\cventry{August 2021 - Present (~5 years)}{AI-Augmented Staff Engineer \& I+D Tech Lead}{Uqomm 
 (\url{www.uqomm.com})}{Con Con, Chile}{}{
Head of the Software and Firmware R\&D department, leading a team of 3 engineers. Responsible for the architectural design, firmware implementation, and deployment of mission-critical communication equipment and telemetry for underground mining.
\begin{itemize}
    \item \textbf{Mission-Critical Embedded Systems (C/C++ Bare-Metal \& FreeRTOS):} Lead firmware architect for underground communication nodes (\textit{fw-vlad}, \textit{fw-gateway2lora}, \textit{fw-netreference}). Implemented hardware Automatic Gain Control (AGC) with dedicated DAC calibration on VLAD25 (tight 26.5 dB gain target), Semtech SX127x (VHF/LoRa) low-level drivers, microsecond-precision timers, and non-blocking Modbus RTU stacks for Cortex-M0/M4 microcontrollers with independent hardware watchdogs (IWDG) and 30-second frequency rollback failovers.
    \item \textbf{Automated Testbenches \& Hardware QA (Python/TDD):} Architected \textit{sw-testbench-orchestrator}, automating RF test instruments (spectrum analyzers, signal generators) to perform unassisted multi-point calibration, sweeping gain and validating hardware RF envelopes; reduced manual lab testing from hours to minutes. Replaced slow cloud CI with deterministic local suites (\textit{tools/verificar.sh}) and pre-push hooks across all firmware repositories.
    \item \textbf{AIOps, Critical Path \& Telemetry (Python/MCP/React):} Engineered custom MCP servers (\textit{jira-mcp}) performing Critical Path Method (CPM) graph analysis on engineering roadmaps via LLMs. Developed full-stack underground telemetry dashboards (\textit{sw-diagnosticoremoto}) with QR-serial device tracing and live equipment remote control.
\end{itemize}
}

\cventry{4 years 2 months / July 2017 - August 2021}{Hardware Engineer}{BlackGPS 
 (\url{www.blackgps.com})}{Santiago, Chile}{}{
As a Design and Hardware Engineer, I was responsible for designing, manufacturing, programming, testing, and commissioning a device to prevent truck theft using GNSS and GSM signal jammers.
\\
In my role as Hardware Support, I configured and supported hardware for fleet control and management, particularly devices from brands like Teltonika, DCT Syrus, and ERM Starlink.
\\
As a BackEnd Developer, I developed new features on the server, processing information from tracking devices using SpringBoot in Java to integrate CANBus data from trucks into the server via cellular.
\\
In my capacity as a Flutter developer, I implemented new functionalities in the mobile app using Flutter, established Bluetooth communication with the tracking device, and improved the user interface.
}

\section{Lumina — Proprietary Deep Tech R\&D \& Quantitative Systems}
Founder and Systems Architect of Lumina, an independent deep-tech initiative focusing on algorithmic execution, decision theory under uncertainty, and high-reliability AI engineering governance.
\\
\cventry{2025 - Present}{Market Microstructure \& Algorithmic Systems}{monteCarlo}{}{}{
Architected a real-time quantitative trading and microstructure research engine in C++20 and Python. Implemented asynchronous WebSockets ingesting 50-level Order Book snapshots, trade flows, and liquidation events with sub-second latency, persisting over 1.6~GB of live market state in SQLite WAL mode. Designed an 8-gate statistical validation framework spanning 690+ days of historical data to eliminate backtest overfitting, coupled with fractional Kelly position sizing and automated walk-forward MLOps pipelines.
}

\cventry{2025 - Present}{Quantitative Decision Engine Under Uncertainty}{desicion-maker}{}{}{
Developed a mathematical decision-support engine integrating Monte Carlo simulations with strict convergence criteria, Value of Information (VoI) calculations, and Multi-Criteria Decision Analysis (TOPSIS supporting negative values and pre-ranking vetoes). Formally validated against academic literature and property-based test suites (\textit{pymcdm}) to evaluate fault-tolerant computing and power-supply hardware architectures.
}

\cventry{2025 - Present}{Enterprise AI Governance Framework}{dev-agents}{}{}{
Architected a master governance and constitution framework for AI coding agents (Claude Code, Cursor, Copilot, Antigravity). Enforces MISRA-C rules, C++20 memory safety invariants, layer-boundary architectures, and local pre-push validation suites, eliminating AI regressions in production embedded codebases.
}

\section{Education}
\cventry{2006--2014}{Electrical Engineer - Computers Major}{Federico Santa María Technical University}{Valparaíso, Chile}{}{}
\section{Thesis}
\cvitem{Title}{\emph{Implementation of a Path Generator in a Navigation System for a Cognitive Mobile Robot}}
\cvitem{Description}{Integrated a route generator into the Cognitive Memory and SLAM (Simultaneous Localization and Mapping) features of the mobile robot IRMA-III to enable autonomous movement in search of a target.}

\section{Technical Skills}
\cvcomputer{\textbf{Programming Languages}}{C, C++, Java, Python}{\textbf{Networking \& Communications}}{TCP/IP, LoRa, GPS, I2C, SPI}
\cvcomputer{\textbf{Scripting}}{Bash, Ansible}{\textbf{Version Control}}{Git}
\cvcomputer{\textbf{Microcontroller Programming}}{STM32, ESP32, Arduino}{\textbf{Web Frameworks}}{React, Express}
\cvcomputer{\textbf{Cloud Platforms}}{Google Cloud Compute Engine}{\textbf{Mobile Development}}{Flutter}
\cvcomputer{\textbf{Operating Systems}}{FreeRTOS, Linux}{\textbf{Tools}}{Open Source}


\section{Management and Interpersonal Skills}
\cvitem{Technical Experience}{Firmware and software development, basic electronic circuit design, technology integration.}
\cvitem{Project Management}{Coordination of cross-functional teams to achieve objectives.}
\cvitem{Innovation}{Participation in the development of new technologies and features.}
\cvitem{User-Centered Approach}{Consideration of user needs and experiences in product development.}
\cvitem{Adaptability}{Ability to quickly learn and adapt to changes in products and industries.}



\section{Languages}
\cvitem{Spanish}{Native}
\cvitem{English}{Advanced Reading, Intermediate Speaking}

\section{Contests and Workshops}

\cventry{November 2022}{HoruS Management Strategy}{Agile Methodologies Training}{Uqomm}{}{}{}
\cventry{July 2019}{Powered by eClass}{Agile Development: Scrum + Kanban}{BlackGps}{}{}{}
\cventry{January 2017}{Summer School}{Analog Integrated Circuit Design}{Synopsys - UTFSM}{}{}{}
\cventry{August 2015}{Winner of I+D Applied Contest with HH Motors project}{Apply Your Idea}{Copec UC Foundation}{}{}
\cventry{October 2015}{Contest-Exhibition of Technological Projects}{EXPOTEC}{UTFSM}{}{}

\end{document}
